\section{RT-X: ¿pueden los Cross-Embodiment Data seguir ampliando la Transfer?}

Si los datos de un solo robot son demasiado pocos, una dirección natural es combinar trajectories de varias instituciones, varios robots y varias tareas. Open X-Embodiment / RT-X unifica en lo posible el action space, la observation y el dataset schema, y luego examina si un modelo puede beneficiarse de distintas embodiment. Es un campo de prueba más cercano a la realidad para la scaling law en robotics.

\subsection{Data mixture y action heterogeneity}

La diapositiva del collage muestra que los datos provienen de distintos brazos robóticos, camera, workspace y tareas. Al combinarlos hay que resolver la action dimension, el coordinate frame, la control frequency, la gripper semantics y los missing sensors. Un simple concatenate puede hacer que un dataset grande opaque a uno pequeño, y también puede hacer que action labels incompatible produzcan negative transfer.

\lecturefigure{slide-37-rtx-data-collage.jpg}{Open X-Embodiment: mezcla de datos entre instituciones y entre robots}{Diapositiva V037 recuperada de la grabación oficial; 00:56:33}

\lecturefigure{slide-38-rtx-positive-transfer.jpg}{RT-X: resultados de positive transfer del entrenamiento entre embodiment}{Diapositiva V038 recuperada de la grabación oficial; 00:57:09}

\begin{knowledgebox}{Términos clave: las cuatro dimensiones del robotics scaling}
\begin{tabularx}{\textwidth}{@{}L{0.18\textwidth}X@{}}
\toprule
Dimensión & Qué se incrementa y riesgo principal \\
\midrule
Model scale & Parámetros y compute; el riesgo es que la data sea insuficiente y que aumenten la latency y la memory \\
Task scale & Variedad de instrucciones y comportamientos; el riesgo es el conflicto de label/action y la long-tail failure \\
Environment scale & Escenarios, objetos y fondos; el riesgo es la observation shift y las affordance nuevas \\
Embodiment scale & Morfología del robot y action space; el riesgo es la incompatibilidad de coordenadas, frecuencias y semántica de control \\
\bottomrule
\end{tabularx}
\end{knowledgebox}

\teachervoice{La síntesis que hace Fei Xia de RT-X es mesurada: en el campo de la robótica se han visto «signs of life» de positive transfer, pero la comunidad de modelos de lenguaje ya considera la transfer algo habitual. Dicho de otro modo, el cross-embodiment scaling es una señal importante, pero todavía no llega a la etapa madura en la que «más datos implica necesariamente mejores resultados».}

\subsection{Conclusiones de la Part 1}

La primera parte concluye en tres puntos: la model consolidation acerca el high-level reasoning y el low-level control a un mismo sistema; el joint scaling pone a prueba si los diverse data producen transfer; y la positive transfer es el criterio decisivo para evaluar si el scaling se sostiene. A continuación, el curso deja de ampliar una sola end-to-end policy y adopta otra interfaz: la salida del LLM pasa a ser reward code.

\lecturefigure{slide-39-part1-summary.jpg}{Part 1: model consolidation, joint scaling y positive transfer}{Diapositiva V039 recuperada de la grabación oficial; 00:57:51}

\subsection{Resumen de la sección}

RT-X convierte el problema de la escasez de datos en un problema de estandarización de datos entre embodiment y de transfer. Extiende la línea de RT-2, pero la action heterogeneity y el dataset imbalance hacen más difícil el «modelo unificado»; esto también explica por qué la segunda línea opta por conservar el controller tradicional.
